
\documentclass[11pt,a4paper]{article}
\usepackage[utf8]{inputenc}
\usepackage[T1]{fontenc}
\usepackage[spanish]{babel}
\usepackage{geometry}
\usepackage{graphicx}
\usepackage{hyperref}
\usepackage{xcolor}
\usepackage{tabularx}
\geometry{margin=2cm}
\setlength{\parindent}{0pt}
\setlength{\parskip}{0.55em}
\sloppy
\definecolor{tcTeal}{HTML}{0E7C74}
\definecolor{tcDark}{HTML}{101817}
\definecolor{tcMuted}{HTML}{526A65}
\definecolor{tcWarnBg}{HTML}{FFF7E5}
\definecolor{tcWarnBorder}{HTML}{C99123}
\hypersetup{colorlinks=true, urlcolor=tcTeal, linkcolor=tcTeal}
\begin{document}

\begin{center}
{\huge\bfseries Informe AI Analyst: MARKET multi market summary}\\[0.25em]
{\large Reporte read-only para revision humana}
\end{center}

\vspace{0.2em}
\hrule height 0.7pt
\vspace{0.9em}

\begin{tabularx}{\textwidth}{>{\bfseries}p{3.2cm}X}
Paquete & market summary e6c3dd63 \\
Setup & MARKET multi market summary \\
Modelo & gpt-5.5 \\
Thinking & medium \\
Macro & not requested \\
Estado & reviewed \\
\end{tabularx}

\vspace{0.75em}
\fcolorbox{tcWarnBorder}{tcWarnBg}{%
\begin{minipage}{0.94\textwidth}
\textbf{Limite de uso.} Este informe no es una senal, no autoriza MT5 y no ejecuta ordenes. Sirve para priorizar revision humana y documentar el contexto.
\end{minipage}}

\vspace{0.4em}

\section*{Resumen}
Revision read-only de mercado multi-timeframe. El conjunto muestra sesgo bajista dominante en marcos cortos: M15 tiene 40 bajistas frente a 20 alcistas y H1 tiene 36 bajistas frente a 24 alcistas. En marcos superiores el balance se modera e incluso gira: H4 presenta 35 alcistas frente a 25 bajistas y D1 32 alcistas frente a 28 bajistas. La lectura principal es de presion correctiva de corto plazo dentro de una estructura diaria menos deteriorada, con varios activos en extremos de RSI o estocastico.

\section*{Lectura tecnica}
El paquete es de estudio, no operativo, sobre 60 instrumentos y con fuentes generadas entre el 6 y el 7 de junio de 2026. La distribucion visual confirma el choque entre cortos plazos bajistas y marcos altos mas equilibrados o alcistas. Hay grupos con alineacion bajista clara, especialmente criptos, metales y parte de commodities, mientras indices y divisas muestran mas dispersion y casos de sobreventa dentro de tendencias superiores alcistas.



\section*{Grafico de referencia}
\begin{center}
\includegraphics[width=0.96\textwidth]{\detokenize{ai_analyst_report_chart.png}}
\end{center}


\section*{Lectura de confluencias}
M15 y H1 concentran mayoria bajista, coherente con lecturas de RSI muy bajas y estocasticos deprimidos en numerosos instrumentos; Criptos como BTCUSD, ETHUSD, LTCUSD y RPLUSD aparecen alineadas a la baja en H1, H4 y D1, con volatilidad claramente por encima de su mediana y rangos 24h amplios; Metales como XAGUSD, XAUEUR, XAUUSD, XPDUSD y XPTUSD muestran predominio bajista multi-timeframe, reforzado por lecturas de impulso debiles; Varios pares vinculados a AUD y NZD muestran sobreventa en H4 dentro de D1 alcista, lo que explica su clasificacion como extremo de RSI en contexto, no como lectura direccional simple; y La volatilidad elevada frente a mediana aparece en 5 casos resumidos, con mayor concentracion en cripto, metales y algunos indices estadounidenses.

\section*{Contradicciones y cautelas}
El mercado agregado no esta alineado de forma limpia: M15/H1 son mayoritariamente bajistas, mientras H4/D1 mantienen mas peso alcista; Hay instrumentos con D1 alcista pero presion bajista en M15-H1-H4, como AUDUSD.r, GBPUSD.r, BRENT y WTI, lo que sugiere correccion contra marco superior; En indices estadounidenses, US100 y US500 aparecen con H4/D1 alcistas pero M15/H1 bajistas y RSI extremo; esta mezcla requiere separar rebote tecnico de continuidad de tendencia; Algunos exoticos USD muestran sobrecompra fuerte en M15/H1/H4 aunque D1 no siempre acompana, especialmente USDMXN.r, USDNOK.r y USDSGD.r; y Existen filas con fechas antiguas frente al resto del paquete, como USDJPY de marzo de 2026 e ITA40 de abril de 2026, que no deben mezclarse sin validar vigencia.

\section*{Riesgos principales}
El paquete es read-only y no contiene permiso operativo; no debe interpretarse como autorizacion de ejecucion; La divergencia entre corto y largo plazo aumenta el riesgo de lecturas falsas si se resume el mercado con una sola direccion; Los activos con RSI extremo pueden permanecer extremos durante fases tendenciales; la sobreventa o sobrecompra no implica por si sola agotamiento; La dispersion temporal de algunas filas introduce riesgo de datos obsoletos en la fotografia agregada; y La volatilidad elevada en cripto, metales y algunos indices puede amplificar ruido intradia y deteriorar comparaciones simples entre instrumentos.

\section*{Contexto macro y noticias}
No se solicito lectura macro para este paquete. El informe se limita al contexto tecnico y a los datos estructurados del Trading Center.

\section*{Conclusion operativa y financiera}
Conclusion operativa y financiera: MARKET queda como market summary de revision, no como instruccion de mercado. La prioridad de revision indicada por el analisis es 2/5 y el riesgo macro figura como no solicitado. La interpretacion conjunta debe apoyarse en tres filtros humanos: que el precio confirme el comportamiento descrito en el grafico, que las contradicciones de tendencia o estructura no invaliden el caso, y que no exista un evento macro cercano capaz de distorsionar el movimiento. Sin esas comprobaciones, el informe solo justifica seguimiento y documentacion, no ejecucion.

\section*{Comprobaciones humanas}
\begin{itemize}
  \item Filtrar o actualizar instrumentos con fechas antiguas antes de usar el resumen como fotografia global.
  \item Separar la revision por grupos: Forex Majors, indices, cripto, commodities y metales, porque el sesgo agregado mezcla comportamientos muy distintos.
  \item Revisar los casos con H4/D1 alcistas y M15/H1 bajistas para distinguir retroceso de ruptura estructural.
  \item Contrastar extremos de RSI con niveles de precio, rango 24h en H1 y rango reciente en H4 antes de extraer conclusiones discrecionales.
  \item Verificar que la muestra de volatilidad elevada no dependa de instrumentos con datos desfasados o sesiones incompletas.
\end{itemize}

\section*{Fuentes}
\textbf{Fuentes locales}\par
\begin{tabularx}{\textwidth}{p{4.2cm}X}
latest\_manifest\_json: C:\textbackslash{}Users\textbackslash{}ralr1\textbackslash{}Desktop\textbackslash{}CD\textbackslash{}TFG\textbackslash{}TFG-Raul\_Rodriguez\textbackslash{}artifacts\textbackslash{}tfg\textbackslash{}trading\_center\_latest\textbackslash{}latest\_manifest.json & Fuente local del paquete de revision \\
market\_radar\_csv: C:\textbackslash{}Users\textbackslash{}ralr1\textbackslash{}Desktop\textbackslash{}CD\textbackslash{}TFG\textbackslash{}TFG-Raul\_Rodriguez\textbackslash{}artifacts\textbackslash{}tfg\textbackslash{}trading\_center\_latest\textbackslash{}market\_radar\textbackslash{}market\_radar.csv & Fuente local del paquete de revision \\
ohlc\_csv: C:\textbackslash{}Users\textbackslash{}ralr1\textbackslash{}Desktop\textbackslash{}CD\textbackslash{}TFG\textbackslash{}TFG-Raul\_Rodriguez\textbackslash{}artifacts\textbackslash{}tfg\textbackslash{}trading\_center\_latest\textbackslash{}ohlc\textbackslash{}ohlc\_mtf.csv & Fuente local del paquete de revision \\
chart.png: distribucion de tendencia M15, H1, H4 y D1 & Fuente local del paquete de revision \\
\end{tabularx}

\end{document}
